\chapter{Continuous Integration with GitHub Actions and Vulnerability Analysis}

This chapter documents the implementation of the Continuous Integration layer of the DevSecOps environment, completing the architecture designed in Chapter~5 and closing the CI/CD cycle begun with ArgoCD in Chapter~7. Up to this point, the GitOps architecture operated with a single repository (\texttt{tfm-gitops}) combining Kubernetes manifests and security policies, managed declaratively by ArgoCD. This chapter materialises the \textit{two-repository pattern} announced in Section~7.1: a second repository (\texttt{tfm-demo-app}) is introduced, containing the application source code and its Dockerfile, and the Continuous Integration pipeline is implemented to build the image, subject it to security analysis, and propagate the change to the GitOps repository for automatic deployment.

A central technical feature of the pipeline is the adoption of a \textit{defence in depth} strategy for vulnerability analysis: instead of a single scanning engine, two complementary engines are implemented. Trivy \cite{trivy_docs} analyses the locally built image, and Grype \cite{grype_docs} analyses an SBOM (\textit{Software Bill of Materials}) \cite{syft_docs} generated by Syft from that same image, consulting different CVE databases. The empirical finding that justifies this strategy---Trivy passed the \texttt{python:3.12-slim} image without alerts while Grype detected a HIGH-severity vulnerability that blocked the pipeline---is one of the most relevant results of this work.

The chapter is organised into seven sections. The first justifies the two-repository split and documents its structure. The second describes the image registry configuration and pipeline credentials. The third documents the demo application and its Dockerfiles. The fourth describes the design of the GitHub Actions workflow and its security decisions. The fifth analyses in depth the dual-engine scanning strategy and the key empirical finding. The sixth covers the technical incidents and lessons learned during implementation. The seventh documents the closure of the complete CI/CD cycle and the integration with the GitOps repository.

% ─────────────────────────────────────────────
\section{Application repository design and two-repository strategy}
% ─────────────────────────────────────────────

The separation between the application repository and the configuration repository is not an organisational preference, but a concrete security requirement justified in Section~5.5.1. The central principle is credential segregation: the CI pipeline needs access to the image registry (\texttt{quay.io}) and write access to the configuration repository, but must under no circumstances have direct access to the Kubernetes cluster. ArgoCD, in turn, needs access to the configuration repository to synchronise the cluster state, but has no reason to access the application source code. This separation establishes a security perimeter between the two phases of the CI/CD cycle, so that the compromise of one does not automatically imply the compromise of the other.

The \texttt{tfm-demo-app} repository contains exclusively the application-layer artefacts:

\begin{itemize}
  \item \texttt{app/}: Python Flask application source code, including the main module (\texttt{app.py}) and the dependency file (\texttt{requirements.txt}).
  \item \texttt{docker/}: the two project Dockerfiles (\texttt{Dockerfile} and \texttt{Dockerfilevulnerable}).
  \item \texttt{.github/workflows/}: the GitHub Actions workflow (\texttt{build-and-push.yaml}) implementing the complete CI pipeline.
\end{itemize}

The \texttt{tfm-gitops} repository, managed by ArgoCD since Chapter~7, is not modified by the development team: only the CI pipeline updates it, and only the \texttt{image} field of the demo application manifest. This delimitation ensures that the security policy manifests (directory \texttt{policies/}) and the ArgoCD \texttt{Application} configuration remain under strict, separate control from the application development cycle.

This architecture materialises the flow described in Section~5.3.1: a commit to the application repository triggers the CI pipeline, which culminates with a commit to the configuration repository, which in turn acts as the trigger signal for ArgoCD. The cluster never receives direct instructions from the pipeline: all interaction with Kubernetes is mediated by the GitOps repository and the reconciliation controller.

% ─────────────────────────────────────────────
\section{Image registry configuration and credentials}
% ─────────────────────────────────────────────

The chosen image registry is Quay.io \cite{quay_docs}, in line with the tools presented in the course and as an enterprise-grade alternative to Docker Hub. Pipeline authentication against the registry is performed via a \textbf{Robot Account}, a central concept in credential management for CI/CD systems that warrants explicit technical justification.

A Robot Account is a non-personal service account whose credentials are decoupled from any individual user. In a CI/CD pipeline, using personal credentials presents two security problems: the first is the \textit{blast radius} in case of compromise---personal credentials have access to all repositories of the user---; the second is operational continuity---if the user leaves the organisation, the pipeline stops working---. A Robot Account addresses both problems: it has access exclusively to the repositories explicitly assigned to it, with the minimum permission level required for its function.

The Robot Account created is named \texttt{zelkarrabi+tfm\_ci\_robot}. Figure~\ref{fig:ci-robotaccount} shows its permission configuration in Quay.io: it has \textit{Write} permission on the \texttt{tfm-demo-app} repository and \textit{None} permission on the \texttt{prova} repository, which is an unrelated test repository. This granularity confirms the application of the least-privilege principle: the Robot Account can only publish images in the specific repository for which it was created, and cannot modify its configuration or access other repositories.

\begin{figure}[H]
\centering
\includegraphics[width=0.90\linewidth]{ci_cd/bloque1_robotaccount.png}
\caption{Configuration of the Robot Account \texttt{zelkarrabi+tfm\_ci\_robot} in Quay.io: \textit{Write} permission on \texttt{tfm-demo-app} and no access on other repositories}
\label{fig:ci-robotaccount}
\end{figure}

The Robot Account credentials are validated beforehand using \texttt{podman login quay.io} from the cluster control node, confirming that the credentials are functional before configuring them in the pipeline. Figure~\ref{fig:ci-podmanlogin} shows the successful authentication with user \texttt{zelkarrabi}.

\begin{figure}[H]
\centering
\includegraphics[width=0.60\linewidth]{ci_cd/bloque1_podmanLogin.png}
\caption{Registry credential verification: successful authentication at \texttt{quay.io} via \texttt{podman login}}
\label{fig:ci-podmanlogin}
\end{figure}
\vspace{-\parskip}

The pipeline requires three repository secrets in GitHub, whose configuration is shown in Figure~\ref{fig:ci-secrets}:

\begin{figure}[H]
\centering
\includegraphics[width=0.90\linewidth]{ci_cd/bloque2_creaccionSecrets.png}
\caption{The three repository secrets configured in GitHub Actions: \texttt{GITOPS\_PAT}, \texttt{QUAY\_PASSWORD} and \texttt{QUAY\_USER}}
\label{fig:ci-secrets}
\end{figure}

\begin{itemize}
  \item \texttt{QUAY\_USER}: Robot Account name (\texttt{zelkarrabi+tfm\_ci\_robot}) for authentication at Quay.io.
  \item \texttt{QUAY\_PASSWORD}: access token generated by Quay.io for the Robot Account.
  \item \texttt{GITOPS\_PAT}: GitHub Personal Access Token for cross-repository writes to \texttt{tfm-gitops}. Unlike the Quay.io Robot Account, a classic PAT with \texttt{repo} scope does not apply the least-privilege principle: it grants read and write access to \textit{all} the user's repositories, not only \texttt{tfm-gitops}. The correct alternative is a \textit{fine-grained PAT} explicitly limited to \texttt{tfm-gitops} with \texttt{Contents: read/write} permission, which replicates for GitHub the same granularity the Robot Account offers in Quay.io; a \textit{deploy key} is another equivalent option. In this lab environment a classic PAT is used for simplicity, with the awareness that in production it should be replaced with one of those alternatives.
\end{itemize}

The explicit use of \texttt{GITOPS\_PAT} instead of the default \texttt{GITHUB\_TOKEN} is a technical constraint of GitHub Actions: the \texttt{GITHUB\_TOKEN} has permissions only over the repository in which the workflow runs (\texttt{tfm-demo-app}), and cannot write to a different repository. The \texttt{GITOPS\_PAT}, being a personal token with explicit scope over \texttt{tfm-gitops}, is the correct mechanism for cross-repository writes.

% ─────────────────────────────────────────────
\section{The demo application and the Dockerfiles}
% ─────────────────────────────────────────────

The demo application is a minimal Python Flask application whose purpose is to validate the complete CI/CD cycle, not to represent a real production workload. It exposes two endpoints on port~8080:

\begin{lstlisting}
from flask import Flask, jsonify
import os, socket

app = Flask(__name__)

@app.route('/')
def index():
    return jsonify({
        "service": "tfm-demo-app",
        "version": "1.0.0",
        "status": "ok",
        "hostname": socket.gethostname(),
        "ci_pipeline": os.getenv("CI_PIPELINE", "manual")
    })

@app.route('/health')
def health():
    return jsonify({"status": "healthy"})

if __name__ == '__main__':
    app.run(host='0.0.0.0', port=8080)
\end{lstlisting}

The root endpoint (\texttt{/}) returns a JSON object that includes the service name, version, pod hostname, and the value of the \texttt{CI\_PIPELINE} environment variable. This last field is what makes it possible to verify empirically, at the end of the chapter, that the code deployed in the cluster corresponds exactly to the commit that triggered the pipeline. The \texttt{/health} endpoint serves as a readiness probe for Kubernetes components.

\subsection{The secure Dockerfile: technical decisions}

The main Dockerfile incorporates several security-oriented design decisions that merit explicit justification:

\begin{lstlisting}
# Stage 1: builder
FROM python:3.13-slim AS builder
WORKDIR /build
COPY app/requirements.txt .
RUN pip install --no-cache-dir --user --upgrade pip && \
    pip install --no-cache-dir --user -r requirements.txt

# Stage 2: runtime
FROM python:3.13-slim
RUN groupadd -r -g 10001 appuser && \
    useradd -r -u 10001 -g appuser appuser

WORKDIR /app
COPY --from=builder --chown=appuser:appuser /root/.local /home/appuser/.local
COPY --chown=appuser:appuser app/app.py /app/

ENV PATH="/home/appuser/.local/bin:${PATH}" \
    PYTHONUNBUFFERED=1 \
    PYTHONDONTWRITEBYTECODE=1

USER appuser
EXPOSE 8080
CMD ["python", "/app/app.py"]
\end{lstlisting}

\textbf{Multi-stage build.} The separation between the \texttt{builder} stage and the runtime stage is the highest-impact decision on the attack surface of the image. The \texttt{builder} stage installs dependencies with \texttt{pip install --user}, which places them in \texttt{/root/.local}. The runtime image copies only that directory---with \texttt{--chown=appuser:appuser}---to the \texttt{/home/appuser/.local} directory of the application user. This way, the \texttt{pip} binary and its own transitive dependencies are not present in the final artefact, reducing the number of scannable components and the exploitation surface at runtime.

\textbf{Minimal base image (\texttt{python:3.13-slim}).} The \texttt{slim} variant of the official Python image removes non-essential operating system packages included in the standard image. This choice represents a balance between compatibility and security: the image includes the system libraries required by most Python applications, but removes development tools, scripting interpreters, and utilities that could be leveraged after an intrusion. The extreme of this strategy is Google's \textit{distroless} \cite{distroless_docs} images, which remove even the shell and package manager, reducing the attack surface to the absolute minimum at the cost of some complexity in runtime failure diagnostics.

\textbf{Non-root user with UID~10001.} The Dockerfile creates \texttt{appuser} as a system account (\texttt{groupadd -r}, \texttt{useradd -r}): without a login shell, without an entry in \texttt{/etc/passwd} with an explicit home, with UID~10001 that avoids collisions with system UIDs (0--999) and typical user UIDs (1000+). The \texttt{USER appuser} instruction makes that identity effective for all container processes. This constitutes a defence-in-depth layer consistent with policy P1 (\texttt{disallow-root}) from Chapter~6: although Kyverno would reject any pod without \texttt{runAsNonRoot:\,true}, the defence starts earlier, in the Dockerfile itself, so that the image operates with minimum privilege regardless of the cluster in which it is deployed.

\textbf{Python environment variables.} \texttt{PYTHONUNBUFFERED=1} ensures that the process standard output is not buffered, allowing logs to be captured in real time by the Kubernetes logging system. \texttt{PYTHONDONTWRITEBYTECODE=1} prevents the writing of \texttt{.pyc} files to the container filesystem, which is consistent with policy P6 (\texttt{add-readonly-rootfs}) that forces the root filesystem to read-only mode: without this variable, Python would attempt to write compiled bytecode and generate permission errors at runtime.

\textbf{Flag \texttt{-{}-no-cache-dir}.} Installing dependencies with this option prevents pip from caching downloaded packages inside the image, reducing its final size and eliminating unnecessary content that might contain outdated package versions.

\subsection{The vulnerable Dockerfile as a reference artefact}

The repository includes a second Dockerfile, \texttt{Dockerfilevulnerable}, based on \texttt{python:3.4-alpine}. Python~3.4 reached end of life in March 2019 and accumulates hundreds of documented CVEs with no fix available in that branch. This artefact was originally designed as a controlled test mechanism to demonstrate the pipeline blocking scenario when facing images with critical vulnerabilities. However, as documented in Section~8.6, this artificial scenario proved unnecessary: the development cycle itself organically produced a genuine block that constitutes empirical evidence of greater academic value than any artificial test.

% ─────────────────────────────────────────────
\section{Continuous Integration pipeline design}
\label{sec:pipeline-design}
% ─────────────────────────────────────────────

The pipeline is defined in the file \texttt{.github/workflows/build-and-push.yaml}. The workflow name is \textit{Build, Scan and Push to Quay} and its single job is named \textit{Build, Trivy scan and push}. The trigger is \texttt{push} on the \texttt{main} branch with path filters that restrict execution to commits affecting the \texttt{app/} directory, the \texttt{docker/} directory, or the workflow file itself:

\begin{lstlisting}
on:
  push:
    branches: [main]
    paths:
      - 'app/**'
      - 'docker/**'
      - '.github/workflows/**'
\end{lstlisting}

This configuration prevents commits affecting only documentation or auxiliary files from unnecessarily triggering the pipeline, reducing GitHub Actions minute consumption and noise in the execution history.

\subsection{Pipeline security decisions}

The four pipeline security design decisions merit explicit technical justification, as they directly materialise the principles stated in Chapter~5.

\textbf{Version pinning via SHA hash.} All third-party GitHub Actions used in the pipeline are referenced by their full commit SHA hash, rather than by semantic version tags or floating tags such as \texttt{@v4} or \texttt{@latest}. For example:

\begin{lstlisting}
- uses: actions/checkout@b4ffde65f46336ab88eb53be808477a3936bae11  # v4.1.1
\end{lstlisting}

This practice materialises the defence against the attack vector documented in Section~5.5.6: a tag like \texttt{@v4} can be redirected by the action maintainer (or by an attacker who has compromised their account) to a different commit without the reference name changing. A SHA hash, by contrast, is immutable: if the workflow manifest references SHA \texttt{b4ffde65f46\ldots}, the runner will execute exactly that commit, regardless of what happens to the repository tags. The limitation of this mechanism, documented as an incident in Section~8.6, is that it does not protect against transitive dependencies of the action code itself.

\textbf{Minimum runner permissions.} The job permission block explicitly declares only the necessary permissions:

\begin{lstlisting}
permissions:
  contents: read
\end{lstlisting}

The default \texttt{GITHUB\_TOKEN} in GitHub Actions has write permissions over the repository. This explicit declaration of \texttt{contents: read} revokes those default permissions, so that even if the runner environment were compromised during execution, the exfiltrated token could not modify the application repository.

\textbf{Image tag based on the commit short SHA.} The tag of the image published in Quay.io is computed from the short SHA of the commit that triggered the pipeline:

\begin{lstlisting}
- name: Compute image tag
  id: vars
  run: |
    SHORT_SHA=$(echo "${GITHUB_SHA}" | cut -c1-7)
    echo "SHORT_SHA=${SHORT_SHA}" >> $GITHUB_OUTPUT
    echo "FULL_IMAGE=${REGISTRY}/${IMAGE_NAME}:${SHORT_SHA}" >> $GITHUB_OUTPUT
\end{lstlisting}

This decision implements the supply chain traceability and immutability principle: each image published in the registry corresponds bijectively to a specific commit in the application repository, making it possible to audit at any time exactly which code is running in the cluster. This design is also consistent with policy P9 (\texttt{disallow-latest-tag}) introduced in Section~7.5.2, which rejects pods whose images are tagged with \texttt{:latest}: the pipeline never produces floating tags, only commit-SHA-based tags.

\textbf{Separation of build and push.} The image is built locally on the runner without being published (\texttt{push: false}). Vulnerability scanning runs against the locally built image. Only if all security controls pass their thresholds is the publication step executed. This sequence implements the \textit{fail fast} principle: a vulnerable image never reaches the registry, eliminating the risk of ArgoCD detecting and deploying it before the incident is identified and managed.

\subsection{Complete pipeline step sequence}

Figure~\ref{fig:ci-workflow-steps-final} shows the step view of the complete workflow in its final version, with the 15 steps that form the pipeline in a successful state:

\begin{figure}[H]
\centering
\includegraphics[width=0.70\linewidth]{ci_cd/bloque5_WorkflowCompleto.png}
\caption{Complete \textit{Build, Trivy scan and push} pipeline: the 15 steps of the CI/CD cycle, executed in 1~min~7~s}
\label{fig:ci-workflow-steps-final}
\end{figure}

The step sequence is as follows:

\begin{enumerate}
  \item \textbf{Set up job}: initialisation of the \texttt{ubuntu-latest} runner and download of SHA-referenced actions.
  \item \textbf{Checkout source}: cloning of the \texttt{tfm-demo-app} repository via SHA-pinned \texttt{actions/checkout}.
  \item \textbf{Compute image tag}: computation of the image tag from the first 7 characters of the commit SHA (\texttt{GITHUB\_SHA | cut -c1-7}), exported as \texttt{SHORT\_SHA} and \texttt{FULL\_IMAGE} outputs.
  \item \textbf{Login to Quay.io}: authentication at \texttt{quay.io} with the Robot Account credentials stored in the \texttt{QUAY\_USER} and \texttt{QUAY\_PASSWORD} secrets.
  \item \textbf{Build Docker image}: image build using \texttt{docker/build-push-action} with \texttt{push:\,false}, producing the image in the runner's local daemon without publishing it.
  \item \textbf{Trivy vulnerability scan}: scanning of the local image with Trivy, configured to fail the pipeline on \texttt{HIGH} or \texttt{CRITICAL} severity vulnerabilities (\texttt{ignore-unfixed: true}).
  \item \textbf{Generate SBOM with Syft}: generation of the Software Bill of Materials in syft-json format (Syft's native format \cite{syft_docs}), producing the file \texttt{sbom.syft.json}.
  \item \textbf{Upload SBOM artifact}: upload of the SBOM as a workflow artefact, accessible from the GitHub Actions interface and downloadable for later auditing.
  \item \textbf{Vulnerability scan with Grype}: SBOM analysis with Grype \cite{grype_docs}, configured to fail on \texttt{HIGH} or higher severity with a fix available (\texttt{--only-fixed}).
  \item \textbf{Push image to Quay}: conditional publication of the image to the registry. Only executes if all previous scanning steps have succeeded.
  \item \textbf{Checkout GitOps repository}: cloning of the \texttt{tfm-gitops} repository using the \texttt{GITOPS\_PAT}.
  \item \textbf{Install yq}: installation of \texttt{yq}, the command-line YAML processor.
  \item \textbf{Update image tag in GitOps repo}: update of the \texttt{image} field in the demo application's \texttt{deployment.yaml} manifest.
  \item \textbf{Commit and push to GitOps repo}: commit and push of the updated manifest to the GitOps repository.
  \item \textbf{Final summary}: execution summary.
\end{enumerate}

The first operational version of the pipeline contained only steps 1--6 and 10 (Trivy as the sole scanning engine, without SBOM, Grype, or the GitOps cycle closure). Figure~\ref{fig:ci-workflow-steps-inicial} documents this initial version, with 50~s execution time, which served as the basis for validating the Quay.io integration before adding the SBOM analysis layer.

\begin{figure}[H]
\centering
\includegraphics[width=0.90\linewidth]{ci_cd/bloque3_workflowSteps.png}
\caption{Initial version of the \textit{Build, Trivy scan and push} pipeline (without SBOM or Grype): build, Trivy scan and publication steps, completed in 50~s}
\label{fig:ci-workflow-steps-inicial}
\end{figure}

The first successful execution of this pipeline published to Quay.io the image with tag \texttt{4888886} (short SHA of the commit), visible in Figure~\ref{fig:ci-primer-tag}. Quay.io reports for this image 8 medium-severity vulnerabilities and 10 fixable ones, with a size of 47.2~MiB and the manifest SHA256 \texttt{3564a6559cbc}.

\begin{figure}[H]
\centering
\includegraphics[width=0.90\linewidth]{ci_cd/bloque3_creaccionTag.png}
\caption{First image published to the \texttt{zelkarrabi/tfm-demo-app} Quay.io repository: tag \texttt{4888886} based on the commit short SHA (47.2~MiB)}
\label{fig:ci-primer-tag}
\end{figure}

% ─────────────────────────────────────────────
\section{Vulnerability analysis: defence in depth with Trivy and Grype}
\label{sec:vuln-analysis}
% ─────────────────────────────────────────────

The incorporation of a second vulnerability analysis engine is not a redundancy decision, but one of technical complementarity. To understand why two engines are more secure than one, it is necessary to understand the structural differences between their analysis methodologies and data sources.

\subsection{Complementary scanning architectures}

\textbf{Trivy} \cite{trivy_docs} analyses the image directly, inspecting its layers and extracting installed packages to cross-reference them against its curated database, composed of sources from Aqua Security and the security data repositories of Linux distributions (Red Hat Security Advisories, Debian Security Tracker, Alpine SecDB, among others). Its strength lies in coverage of operating system packages and the speed of its database, which is updated at high frequency.

\textbf{Syft + Grype} operates in two phases. Syft \cite{syft_docs} generates an SBOM of the image, which is a structured and standardised inventory of all artefact components (operating system packages, language libraries, binaries). Grype \cite{grype_docs} analyses that SBOM against its own database, which aggregates sources different from Trivy's: Anchore Enterprise Data, GitHub Security Advisories, National Vulnerability Database (NVD), and Python-specific sources such as the PyPI Advisory Database. This difference in data sources means Grype can detect vulnerabilities that Trivy has not yet incorporated into its database, and vice versa.

\subsection{The SBOM as a compliance deliverable}

Generating the SBOM with Syft is not merely an intermediate step for Grype analysis: the produced file is a \textit{compliance} and traceability artefact with inherent value. An SBOM documents, in a structured and standardised format, which components form part of each artefact, with their exact versions and licences. This inventory allows audit questions to be answered such as: what version of \texttt{urllib3} is in production? which artefacts contain a library affected by a newly published CVE?

There are two main standards for SBOM specification:

\begin{itemize}
  \item \textbf{SPDX} \cite{spdx_spec}: the Linux Foundation standard, primarily oriented toward software licence management and compliance with NIST SSDF \cite{nist_ssdf} software supply chain requirements.
  \item \textbf{CycloneDX} \cite{cyclonedx_spec}: the OWASP Foundation standard, specifically oriented toward security and the DevSecOps lifecycle, with native support for vulnerability analysis and risk metadata.
\end{itemize}

The pipeline emits the SBOM in \textbf{syft-json} format, Syft's native format, which Grype consumes directly for vulnerability analysis. The generated artefact (\texttt{sbom.syft.json}) is uploaded to the workflow under the name \texttt{sbom-\$\{SHORT\_SHA\}}, associating it with the specific run and corresponding image. Unlike SPDX or CycloneDX, syft-json is not a recognised interoperability standard but the tool's internal format; in an environment with formal compliance requirements, the pipeline could additionally emit the SBOM in CycloneDX via \texttt{-o cyclonedx-json}. Figure~\ref{fig:ci-sbom-upload} shows the successful execution of the \textit{Upload SBOM artifact} step: the file \texttt{sbom-defe4ef.zip} was uploaded with a final size of 490,279~bytes (artefact ID \texttt{7196257107}), accessible in the workflow artefacts interface \cite{github_actions_docs}.

\begin{figure}[H]
\centering
\includegraphics[width=0.95\linewidth]{ci_cd/bloque4_ArtefactoSBOM.png}
\caption{\textit{Upload SBOM artifact} step: upload of \texttt{sbom-defe4ef.zip} as an audit artefact (490,279~bytes, ID~\texttt{7196257107})}
\label{fig:ci-sbom-upload}
\end{figure}

\subsection{Empirical finding: CVE undetected by Trivy, blocking for Grype}

The most relevant empirical finding of the chapter occurred during the first execution of the pipeline with the dual scanning engine active, using the base image \texttt{python:3.12-slim}. Trivy completed its scan without reporting vulnerabilities above the configured severity threshold, and the step advanced correctly. Grype, however, analysed the SBOM generated by Syft and detected vulnerability \textbf{CVE-2025-13836} with \textbf{HIGH} severity in the Python~3.12.13 binary, with fix versions available from \texttt{3.13.11}, \texttt{3.14.1} and \texttt{3.15.0}. The pipeline failed with exit code~1 in the \textit{Vulnerability scan with Grype} step, preventing image publication.

Figure~\ref{fig:ci-grype-output} shows the complete output of the failed step: in addition to CVE-2025-13836 (HIGH), Grype reported multiple MEDIUM-severity vulnerabilities in the Python~3.12.13 binary (CVE-2026-2297, CVE-2026-1299, CVE-2026-0865, CVE-2025-6075, among others) and in the installed Python dependencies (\texttt{flask~3.0.3}, \texttt{pip~25.0.1}, \texttt{werkzeug~3.0.4}). The log closing line, \texttt{Error: Process completed with exit code 1}, confirms that the pipeline stopped at this point, without executing the publication or GitOps repository update steps.

\begin{figure}[H]
\centering
\includegraphics[width=0.95\linewidth]{ci_cd/bloque4_outputTrivyHighPython.png}
\caption{Output of the \textit{Vulnerability scan with Grype} step (\texttt{anchore/grype:v0.79.6}): pipeline blocked by CVE-2025-13836 (HIGH) in \texttt{python~3.12.13}, with fix versions from \texttt{3.13.11}}
\label{fig:ci-grype-output}
\end{figure}

This behaviour is not a failure of either tool. Trivy and Grype have different databases, different severity taxonomies, and independent rates of incorporating new CVEs. CVE-2025-13836 was present in Grype's database at the time of execution but not in Trivy's. This temporal lag between databases is a known and documented phenomenon in the container security literature \cite{nist_sp800190}: organisations that rely on a single scanning engine implicitly assume that engine has complete coverage of all relevant CVEs, an assumption this experiment refutes empirically.

The conclusion is precisely the inverse: diversity of analysis engines provides real security layers because each engine can detect CVEs that others have not yet incorporated. The defence-in-depth principle, typically applied at the cluster access control and configuration layer (Kyverno as the admission engine, ArgoCD as the declarative state controller), is applied here to the image supply chain layer: no scanning engine is sufficient on its own.

It is worth, however, bounding the scope of the finding. CI scanning is \textit{point-in-time}: it only detects vulnerabilities present in the databases at the time of the build; a CVE published after the image has been built and published will go unnoticed until the next cycle. This makes periodic re-scanning of already published images and runtime monitoring indispensable---a function that in the Chapter~5 architecture corresponds to Sysdig. Additionally, the defence-in-depth conclusion rests on a single empirical case in which the engines employed different severity taxonomies; in scenarios where both taxonomies are aligned, the coverage difference between engines may be smaller.

% ─────────────────────────────────────────────
\section{Technical incidents and lessons learned}
% ─────────────────────────────────────────────

This section documents the three most relevant technical incidents that arose during pipeline implementation. Like the lessons-learned sections of Chapters~6 and~7, the objective is not to narrate failures, but to extract technical conclusions of value that illustrate general principles of secure CI/CD system design.

\subsection{Incident 1: supply chain attack on the Trivy action itself}

The first incident affected the very defence mechanism designed to protect the pipeline from supply chain attacks: SHA-based version pinning of actions. Despite referencing \texttt{aquasecurity/trivy-action} by its full SHA hash, the pipeline failed on the first execution attempt with an action resolution error.

The root cause, identified after analysing the runner logs, was that \texttt{aquasecurity/trivy-action} has an internal transitive dependency on an auxiliary action named \texttt{setup-trivy}, referenced via a mutable tag (not SHA-pinned). Aqua Security had removed or relocated this tag as part of a response to a security incident in their own infrastructure, which broke the action resolution chain even though the SHA of the main action remained intact.

Figure~\ref{fig:ci-workflow-inicial} shows the history of the first two workflow executions: the first (\textit{feat: initial application code, Dockerfiles and CI workflow}, commit \texttt{6070383}) failed in 7~s for this reason; the second (\textit{ci: update trivy-action SHA to fix broken transitive dependency}, commit \texttt{4888886}) succeeded in 52~s after updating to the SHA corresponding to the stable \texttt{v0.35.0} version.

\begin{figure}[H]
\centering
\includegraphics[width=0.95\linewidth]{ci_cd/bloque3_workflowActions.png}
\caption{Initial workflow history: failure of run~\#1 (7~s) due to a broken transitive dependency in the Trivy action, and correction in run~\#2 (52~s) with the \texttt{v0.35.0} SHA}
\label{fig:ci-workflow-inicial}
\end{figure}

This incident empirically demonstrates a real limitation of supply chain security in CI systems based on third-party components: SHA pinning protects only against mutation of the directly referenced action code, but not against transitive dependencies that the action may resolve internally via mutable tags. In a production environment, complete mitigation would require either that action providers also pin their internal dependencies by SHA, or the use of a \textit{vendoring} system that internalises all dependencies in the corporate repository. The pragmatic solution adopted---updating to the SHA of the stable version with active support---is an acceptable compromise that restores pipeline immutability for the most recent version of the action.

\subsection{Incident 2: reproducibility versus Grype database freshness}

The second incident affected the \textit{Vulnerability scan with Grype} step after the SBOM was incorporated. By pinning the Grype image version as \texttt{anchore/grype:v0.79.6}, the vulnerability database included in that image was static and dated from the image's publication date. By default, Grype validates that its CVE database is no more than 5~days old; if the threshold is exceeded, Grype fails with an error before processing the SBOM.

This behaviour exposes a fundamental \textit{trade-off} in security pipeline design: \textbf{reproducibility} (pinned image version, deterministic pipeline behaviour) conflicts with \textbf{database currency} (coverage of the most recent vulnerabilities). Prioritising reproducibility means the pipeline will always use the same version of Grype with the same database, but that database becomes obsolete over time. Prioritising database currency means the Grype image updates automatically, but introduces variability in pipeline behaviour (a run may fail not because the image has new vulnerabilities, but because the updated database incorporates new entries).

The adopted solution was to configure the environment variable \texttt{GRYPE\_DB\_VALIDATE\_AGE=false}, which instructs Grype to skip age validation and use the database available in the pinned image:

\begin{lstlisting}
- name: Vulnerability scan with Grype
  run: |
    docker run --rm \
      -e GRYPE_DB_VALIDATE_AGE=false \
      -v "$PWD":/work \
      anchore/grype:v0.79.6 \
      sbom:/work/sbom.syft.json \
      --fail-on high \
      --only-fixed
\end{lstlisting}

This decision prioritises pipeline determinism: the same workflow with the same inputs will always produce the same behaviour, facilitating incident debugging and result reproduction. In a production environment, the recommended alternative would be a controlled periodic renewal of the Grype image version through a managed update process (weekly automatic pull request, for example), combining database updates with an explicit review before being incorporated into the pipeline.

Figure~\ref{fig:ci-workflow-sbom-iteraciones} documents the complete workflow execution history during the resolution of this incident: runs \#3, \#4 and \#5 correspond to the attempts to resolve the Grype database problem, until reaching the stable solution with \texttt{GRYPE\_DB\_VALIDATE\_AGE=false} in run \#6.

\begin{figure}[H]
\centering
\includegraphics[width=0.95\linewidth]{ci_cd/bloque4_flujoWorkflowAlIncluirSBOMGrype.png}
\caption{Workflow \textit{Build, Scan and Push to Quay} history: incident resolution iterations with SBOM and Grype (runs \#3--\#5 failed) and final successful execution with \texttt{python:3.13-slim} (run \#6, commit \texttt{defe4ef})}
\label{fig:ci-workflow-sbom-iteraciones}
\end{figure}

\subsection{Incident 3: applying the Shift-Left principle against CVE-2025-13836}

The third incident is, in terms of academic value, the most representative of the chapter: the pipeline blocked the publication of the \texttt{python:3.12-slim} image due to CVE-2025-13836 (HIGH) detected by Grype, as documented in Section~\ref{sec:vuln-analysis}. Faced with this block, there were two response paths:

\begin{enumerate}
  \item Reduce the Grype severity threshold from \texttt{HIGH} to \texttt{CRITICAL}, so that CVE-2025-13836 would not block the pipeline.
  \item Fix the vulnerability at the source by updating the base image to a Python version that does not contain the CVE.
\end{enumerate}

The first option was evaluated and rejected because it would compromise the system's security standards: a HIGH-severity vulnerability in the Python runtime is not a false positive or an acceptable risk in a DevSecOps environment, and relaxing the detection threshold would be equivalent to systematically ignoring every vulnerability that does not reach maximum criticality. The second option, updating the base image to \texttt{python:3.13-slim}, resolves the problem at the source without compromising any control.

This decision exemplifies the \textit{Shift-Left} principle \cite{sharma2020devsecops}: instead of managing the vulnerability at the deployment layer (adding exceptions or relaxing controls), it is corrected at the earliest possible phase of the lifecycle---the selection of the Dockerfile base image. The pipeline is not just a vulnerability detector; it is a mechanism that imposes an obligation to fix at the source before the artefact advances through the chain.

The fix was applied in commit \texttt{defe4ef} (\textit{fix: bump base image to python:3.13-slim to remediate High vulnerability}). The resulting image, published to Quay.io under the same tag, has 7 medium-severity vulnerabilities and 8 fixable ones with a size of 46.9~MiB---a reduction in both the number of CVEs and the image size, as expected when using a more modern and patched version of Python.

\textbf{Note on the Dockerfilevulnerable.} The blocking scenario that \texttt{Dockerfilevulnerable} (based on \texttt{python:\allowbreak{}3.4-alpine}) would have allowed reproducing artificially proved unnecessary: the development pipeline itself produced a verified genuine test by legitimately blocking an image the team considered secure at the time (\texttt{python:3.12-slim}). This real-world test has greater academic value than an artificial one because it demonstrates the system working in a production scenario where vulnerabilities emerge without warning: the developer did not deliberately choose an insecure image, nor was the block prepared. The CVE was present in Grype's database, the image contained it, and the pipeline did exactly what it was designed to do.

% ─────────────────────────────────────────────
\section{Closing the CI/CD cycle: integration with the GitOps repository}
% ─────────────────────────────────────────────

The final block of the pipeline closes the CI/CD cycle, connecting the Continuous Integration layer documented in this chapter with the Continuous Deployment layer implemented in Chapter~7. Once all security controls have been passed and the image published to Quay.io, the pipeline updates the GitOps repository automatically.

\subsection{Automatic GitOps repository update}

The update process is performed in four consecutive workflow steps:

\textbf{Checkout GitOps repository.} The pipeline checks out the \texttt{tfm-gitops} repository using the \texttt{GITOPS\_PAT}:

\begin{lstlisting}
- name: Checkout GitOps repository
  uses: actions/checkout@b4ffde65f46336ab88eb53be808477a3936bae11  # v4.1.1
  with:
    repository: ezixd23/tfm-gitops
    token: ${{ secrets.GITOPS_PAT }}
    path: gitops-repo
    ref: main
\end{lstlisting}

Using \texttt{GITOPS\_PAT} instead of the default \texttt{GITHUB\_TOKEN} is a technically justified decision: the \texttt{GITHUB\_TOKEN} has permissions only over the repository in which the workflow runs (\texttt{tfm-demo-app}) and cannot write to a different repository. The \texttt{GITOPS\_PAT} is the necessary mechanism for cross-repository writes; in production it should be a fine-grained PAT limited to \texttt{tfm-gitops} with \texttt{Contents: read/write}, as reasoned in Section~8.2.

\textbf{Install yq.} \texttt{yq}, Mike Farah's command-line YAML processor, is installed to allow modifying specific fields of a YAML file without altering the rest of the file's structure.

The choice of \texttt{yq} over alternatives such as \texttt{sed} is based on robustness considerations. \texttt{sed} operates on plain text without understanding YAML structure: a format change in the file (different spacing, comments, field reordering) can silently break the substitution regular expression, producing a malformed manifest. \texttt{yq}, by processing YAML as a structured data tree, guarantees that only the \texttt{image} field of the first container of the Deployment is modified, preserving the rest of the file unchanged.

\textbf{Update image tag in GitOps repo.} The manifest update is performed with a \texttt{yq} expression that navigates the Deployment structure:

\begin{lstlisting}
- name: Update image tag in GitOps repo
  run: |
    cd gitops-repo
    yq eval -i \
      ".spec.template.spec.containers[0].image = \"${{ steps.vars.outputs.FULL_IMAGE }}\"" \
      applications/demo-app/deployment.yaml
\end{lstlisting}

\textbf{Commit and push to GitOps repo.} The commit includes a message with cross-repository traceability:

\begin{lstlisting}
- name: Commit and push to GitOps repo
  run: |
    cd gitops-repo
    git config user.name "GitHub Actions CI"
    git config user.email "actions@github.com"
    git add applications/demo-app/deployment.yaml
    git commit -m "chore: update demo-app image to ${{ steps.vars.outputs.SHORT_SHA }}"
    git push origin main
\end{lstlisting}

The commit message includes the image short SHA (\texttt{steps.vars.outputs.SHORT\_SHA}), enabling bidirectional auditing: given a commit in \texttt{tfm-gitops}, it is possible to identify the deployed image tag, and given an image tag in Quay.io, it is possible to locate the source code commit that originated it.

\subsection{Empirical evidence of the complete cycle}

Figure~\ref{fig:ci-git-log} shows the commit history in the \texttt{tfm-gitops} repository after the complete pipeline execution. The most recent commit (\texttt{f762c09}, \texttt{HEAD}) has the message \texttt{chore: update demo-app image to d8ae0bd}, direct evidence that the pipeline automatically updated the GitOps repository with the new image tag.

\begin{figure}[H]
\centering
\includegraphics[width=0.90\linewidth]{ci_cd/bloque5_logDeCommitAutomatico.png}
\caption{Commit history in \texttt{tfm-gitops}: the automatic commit \texttt{f762c09} updates the image tag to \texttt{d8ae0bd}, triggering ArgoCD synchronisation}
\label{fig:ci-git-log}
\end{figure}

Once this new commit is detected by ArgoCD through its repository polling mechanism (and its continuous watch over synchronised cluster resources described in Section~7.5.1), the reconciliation controller initiates synchronisation of the \texttt{Deployment} with the new image specification. Figure~\ref{fig:ci-quay-tags} shows the history and status of the image repository in Quay.io, where four tags published by different pipeline runs are now visible: \texttt{d8ae0bd} as the most recent image (46.9~MiB), alongside previous versions \texttt{036fe39}, \texttt{defe4ef} (image with \texttt{python:3.13-slim}, 46.9~MiB) and \texttt{4888886} (initial image with \texttt{python:3.12-slim}, 47.2~MiB).

\begin{figure}[H]
\centering
\includegraphics[width=0.95\linewidth]{ci_cd/bloque4_nuevoTAG.png}
\caption{Repository \texttt{zelkarrabi/tfm-demo-app} on Quay.io showing the history of the four tags published by the pipeline, with \texttt{d8ae0bd} as the most recent and secure image.}
\label{fig:ci-quay-tags}
\end{figure}

Figure~\ref{fig:ci-pod-image} shows the final verification from the cluster: the two pods of the \texttt{demo-app} Deployment in \texttt{Running} state are executing the image \texttt{quay.io/zelkarrabi/tfm-demo-app:d8ae0bd}, the tag based on the short SHA of the last commit that passed all pipeline controls.

\begin{figure}[H]
\centering
\includegraphics[width=0.90\linewidth]{ci_cd/bloque5_podConImagenActualizada.png}
\caption{Cluster verification: the two \texttt{demo-app} pods in \texttt{Running} state with image \texttt{quay.io/zelkarrabi/\allowbreak{}tfm-demo-app:d8ae0bd} deployed by ArgoCD}
\label{fig:ci-pod-image}
\end{figure}

The ArgoCD interface pod view, shown in Figure~\ref{fig:ci-argocd-pod}, presents the full detail of pod \texttt{demo-app-86d6c6d8f7-8s9c6} created on 06/03/2026 at 13:45:41: the active image is \texttt{quay.io/zelkarrabi/tfm-demo-app@sha256:6279046da2c79968\ldots}, the container state is \texttt{Running} (\textit{started and ready}) and the pod health is \texttt{Healthy}. The SHA256 digest reference appearing in the pod's \texttt{imageID} field is not resolved by ArgoCD: it is the kubelet that, at pull time, downloads the tag-referenced image and records the exact digest of the downloaded manifest; ArgoCD merely reflects that live value from the pod state. The practical result is effective immutability: even if the tag were redirected in the registry to a different image, the running pod has already fixed the digest in its \texttt{imageID}.

\begin{figure}[H]
\centering
\includegraphics[width=0.85\linewidth]{ci_cd/bloque5_actualizadoImagenQuaypng.png}
\caption{Detail of pod \texttt{demo-app-86d6c6d8f7-8s9c6} in ArgoCD: active image referenced by SHA256 digest, \texttt{Running} state and \texttt{Healthy} health}
\label{fig:ci-argocd-pod}
\end{figure}

The final end-to-end test consists of verifying that the change introduced in the \texttt{tfm-demo-app} source code is visible in the HTTP response of the application deployed in the cluster. Figure~\ref{fig:ci-curl-end-to-end} shows the application response accessed at \texttt{192.168.1.101:8888}: the returned JSON object includes the field \texttt{ci\_pipeline:\,"end-to-end-test"}, the value of the \texttt{CI\_PIPELINE} environment variable configured in the Deployment, which confirms that the deployed code corresponds exactly to the commit that activated the pipeline. The field \texttt{hostname:\,demo-app-7fdd54f645-4d5nk} identifies the specific pod that served the request, and \texttt{service:\,tfm-demo-app} and \texttt{version:\,1.0.0} correspond to the values declared in the application source code.

\begin{figure}[H]
\centering
\includegraphics[width=0.80\linewidth]{ci_cd/bloque5_curlMostrandoCambioCodigoFuente.png}
\caption{HTTP response from the application at \texttt{192.168.1.101:8888}: the \texttt{ci\_pipeline:\,end-to-end-test} field confirms the complete propagation of the source code change to the runtime environment}
\label{fig:ci-curl-end-to-end}
\end{figure}

\subsection{The complete DevSecOps cycle}

The combination of the results documented in this chapter and Chapter~7 constitutes the complete materialisation of the DevSecOps architecture designed in Chapter~5. The end-to-end flow empirically verified is as follows:

\begin{enumerate}
  \item A developer makes a change to the \texttt{tfm-demo-app} source code and pushes to the \texttt{main} branch.
  \item GitHub Actions detects the event and launches the \textit{Build, Scan and Push to Quay} workflow.
  \item The pipeline builds the Docker image with the multi-stage Dockerfile and non-root user.
  \item Trivy scans the image; if any \texttt{HIGH} or \texttt{CRITICAL} vulnerability with a fix available is detected, the pipeline fails here.
  \item Syft generates the image SBOM and uploads it as a permanent audit artefact.
  \item Grype analyses the SBOM; if any HIGH or higher severity vulnerability is detected, the pipeline fails here.
  \item The image, free of vulnerabilities above the threshold, is published to Quay.io under the commit short SHA tag.
  \item The pipeline updates the \texttt{image} field of the Deployment in \texttt{tfm-gitops} via \texttt{yq} and commits with cross-repository traceability.
  \item ArgoCD detects the new commit in \texttt{tfm-gitops} through its watch on synchronised resources and its Git polling cycle, and synchronises the Deployment with the new image.
  \item Kubernetes creates the new pods; Kyverno intercepts the creation request, applies the mutation policies (P4--P7) and validates against the validation policies (P1--P3, P9). If the pod does not satisfy any condition not fixable through mutation, it is rejected.
  \item The new pods transition to \texttt{Running} state and the application responds HTTP with the values from the new source code.
\end{enumerate}

This flow integrates the three layers of the DevSecOps architecture from Chapter~5: the development layer (source code and Git), the automation layer (GitHub Actions and ArgoCD), and the execution layer (Kubernetes with Kyverno). Security is integrated at every transition between layers: Trivy and Grype between the development layer and the image registry; Kyverno policies between the registry and the cluster; and ArgoCD as the guarantor that the cluster state remains aligned with the configuration repository.

With the complete CI/CD cycle implemented and empirically verified, the next chapter addresses the practical comparison between Kyverno and OPA Gatekeeper, which complements the theoretical review of Chapter~2 with evidence obtained through the direct reimplementation of a representative subset of the Chapter~6 policies in OPA's Rego language.
